\section{Mengakses Elemen DOM dan Mengubah Konten}

JavaScript menyediakan beberapa metode untuk mengakses elemen di DOM. Metode modern yang paling direkomendasikan adalah \code{querySelector()} (mengembalikan elemen pertama yang cocok) dan \code{querySelectorAll()} (mengembalikan semua elemen yang cocok). Keduanya menggunakan CSS selector sebagai argumen, sehingga sangat fleksibel \cite{mdnJsGuide}.

\begin{table}[htbp]
\centering
\begin{tabular}{|P{5cm}|P{2.5cm}|P{4.5cm}|}
\hline
\textbf{Metode} & \textbf{Hasil} & \textbf{Catatan} \\
\hline
\code{getElementById("id")} & 1 elemen & Cepat; hanya by id \\
\hline
\code{querySelector(".css")} & 1 elemen & CSS selector; fleksibel \\
\hline
\code{querySelectorAll("p")} & NodeList & Semua yang cocok \\
\hline
\code{getElementsByClassName()} & HTMLCollection & Live collection; legacy \\
\hline
\end{tabular}
\caption{Metode akses elemen DOM}
\end{table}

Setelah mendapatkan elemen, konten dapat diubah: \code{textContent} mengubah teks (aman dari XSS), \code{innerHTML} mengubah HTML di dalam elemen (hati-hati: rentan XSS jika menyisipkan input pengguna tanpa sanitasi). Atribut diubah dengan \code{setAttribute()} atau properti langsung (\code{el.src = "..."}). Style diubah dengan \code{el.style.color = "red"} atau lebih baik menggunakan class: \code{el.classList.add("active")} \cite{jsInfo}.

\begin{contoh}
\textbf{Studi Kasus: Counter Interaktif}

Tombol + dan - mengubah angka yang ditampilkan di halaman:
\begin{enumerate}
  \item HTML: \code{<span id="count">0</span>} dan dua \code{<button>}
  \item JavaScript: variabel \code{let count = 0;} dan dua event listener
  \item Klik +: \code{count++; document.getElementById("count").textContent = count;}
  \item Klik -: \code{count--; ...} (1 baris)
\end{enumerate}
\end{contoh}

\begin{catatan}
\textbf{textContent vs innerHTML.} \code{textContent} mengubah teks polos---aman karena HTML tag akan di-escape (ditampilkan sebagai teks, bukan dirender). \code{innerHTML} menyisipkan HTML yang akan dirender oleh browser---\textbf{jangan} gunakan untuk menyisipkan input pengguna tanpa sanitasi (risiko XSS). Untuk menambah elemen secara aman, gunakan \code{createElement()} + \code{appendChild()} \cite{owaspXss}.
\end{catatan}

